% year: תשפ"ד
% semester: א
% moed: א
% number: 2
% points: 16
% categories: continuity, func-limits
% source: exams/מבחנים/תשפד/מועד א תשפד.pdf

\begin{question}
יהיו $a,b\in\R$. נגדיר פונקציה על ידי
\[ f(x)=\begin{cases} \frac{\sin(8x)}{x} & x<0,\\[2pt] ax+b & 0\le x\le 4,\\[2pt] \frac{x^2-3x-4}{2-\sqrt{x}} & x>4 \end{cases} \]
לאילו ערכים של $a,b$ הפונקציה רציפה ב-$\R$? נמקו.
\end{question}

\begin{hint}
בכל אחד משלושת התחומים הפתוחים הפונקציה רציפה כהרכבה/מנה של פונקציות רציפות. נותר לבדוק את נקודות התפר $x=0$ ו-$x=4$.
\end{hint}

\begin{hint}
בגבול ב-$x=4$ פרקו את המונה $x^2-3x-4=(x-4)(x+1)$ וכתבו $x-4=(\sqrt x-2)(\sqrt x+2)$.
\end{hint}

\begin{solution}
\textbf{רציפות בתוך כל תחום.}
\begin{itemize}
  \item בתחום $x<0$: $f(x)=\frac{\sin 8x}{x}$ היא מנה של פונקציות רציפות שהמכנה שלה אינו מתאפס, ולכן רציפה.
  \item בתחום $0<x<4$: $f(x)=ax+b$ פולינום, רציף.
  \item בתחום $x>4$: $f(x)=\frac{x^2-3x-4}{2-\sqrt x}$ מנה של פונקציות רציפות, והמכנה מתאפס רק ב-$x=4$ שאינו בתחום. לכן רציפה.
\end{itemize}
לכן $f$ רציפה ב-$\R$ אם ורק אם היא רציפה בנקודות $x=0$ ו-$x=4$.

\textbf{הנקודה $x=0$.} $f(0)=b$, ומימין $\lim_{x\to0^+}f(x)=\lim_{x\to0^+}(ax+b)=b$. משמאל, לפי הגבול היסודי $\lim_{t\to0}\frac{\sin t}{t}=1$:
\[ \lim_{x\to0^-}\frac{\sin(8x)}{x} = \lim_{x\to0^-}8\cdot\frac{\sin(8x)}{8x} = 8. \]
לכן רציפות ב-$0$ מתקיימת אם ורק אם $b=8$.

\textbf{הנקודה $x=4$.} $f(4)=4a+b$, ומשמאל $\lim_{x\to4^-}(ax+b)=4a+b$. מימין, עבור $x>4$:
\[ \frac{x^2-3x-4}{2-\sqrt x} = \frac{(x-4)(x+1)}{2-\sqrt x} = \frac{(\sqrt x-2)(\sqrt x+2)(x+1)}{-(\sqrt x-2)} = -(\sqrt x+2)(x+1), \]
ולכן
\[ \lim_{x\to4^+}f(x) = -(2+2)(4+1) = -20. \]
רציפות ב-$4$ מתקיימת אם ורק אם $4a+b=-20$.

\textbf{מסקנה.} עם $b=8$ נקבל $4a=-28$, כלומר
\[ \boxed{a=-7,\quad b=8}, \]
ורק עבור ערכים אלה $f$ רציפה בכל $\R$.
\end{solution}
